\section{Konzeption der Informationsverarbeitungsarchitektur}
\label{sec:architecture}

Dieses Kapitel entwickelt das Architekturkonzept für die kontrollierte
KI-gestützte Transformation heterogener Engineering-Information in
SysML~v2. Ausgehend von den Anforderungen an Informationsgrundlage,
Interpretation und menschliche Autorisierung werden die Zielmodellbildung,
die regelgebundene Transformation sowie die durchgängige Provenienz und
Nachverfolgbarkeit schrittweise hergeleitet. Anschließend wird der Ansatz auf
die Verarbeitung mehrerer Quellen erweitert und die entwickelten Mechanismen
zu einem integrierten Architekturkonzept zusammengeführt.

\subsection{Architektur- und Gestaltungsziel}
\label{sec:architecture_problem}

Large Language Models ermöglichen die KI-gestützte Erzeugung von
Engineering-Artefakten aus natürlichsprachlicher Information. Untersuchungen
im Systems Engineering zeigen, dass dabei Artefakte entstehen können, die
menschlich erstellten Ergebnissen in ihrer sprachlichen und semantischen
Erscheinung teilweise stark ähneln \cite{Topcu2025}. Auch die Erzeugung
textueller SysML-v2-Artefakte durch LLMs ist grundsätzlich möglich
\cite{Stein2026}. Allgemeine LLMs verfügen jedoch ohne zusätzlich
bereitgestellten Kontext nur begrenzt über projektspezifische
Detailinformationen. Die erzeugten Ergebnisse können daher textuell plausibel
erscheinen und eine hohe Ähnlichkeit zu Expertenlösungen aufweisen, zugleich
aber fachliche Defizite, unbelegte Annahmen oder Überspezifikationen enthalten
\cite{Topcu2025}. Deshalb müssen bei der SysML-v2-Generierung erzeugte
Artefakte zusätzlich fachlich geprüft und bewertet werden
\cite{Stein2026}.

Für den Transformationsprozess ergibt sich daraus die Anforderung, die
fachlichen und interpretationsbasierten Entscheidungen zwischen
Ausgangsinformation und generiertem Modell nachvollziehbar abzubilden. Damit
auch die Entstehung möglicher Fehler oder Abweichungen rekonstruiert werden
kann, müssen die für die Herleitung relevanten Informationszustände,
Verarbeitungsschritte und Entscheidungsgrundlagen nachvollziehbar erhalten
bleiben. Dies entspricht dem grundlegenden Provenienzgedanken,
Informationsobjekte, Verarbeitungsaktivitäten und beteiligte Akteure in ihren
Entstehungs- und Ableitungsbeziehungen miteinander zu verbinden
\cite{MoreauMissier2013}.

Die textuelle und maschinenverarbeitbare Repräsentation von SysML~v2 bietet
eine geeignete Grundlage für einen \emph{Architecture as Code-Ansatz}.
Maschinenlesbare und versionskontrollierte Architekturartefakte unterstützen
unter anderem automatisierte Verarbeitung, Versionierung und
Nachverfolgbarkeit \cite{Bucaioni2025,Stirbu2022}. Das in dieser Arbeit
vorgestellte Architekturkonzept nutzt diese Eigenschaften, um das generierte
Modell nachvollziehbar mit den für seine Herleitung relevanten Quellen,
Verarbeitungszuständen und Entscheidungen zu verknüpfen.

Das Transformationsergebnis umfasst damit neben der formalen
SysML-v2-Repräsentation auch die für ihre Herleitung erforderlichen
Nachverfolgbarkeitsbeziehungen. Sie machen sichtbar, welche Informationen
fachliche Aussagen begründet haben, an welchen Stellen probabilistische
Interpretation stattgefunden hat, welche Ergebnisse zunächst nur
maschinelle Vorschläge oder quellengebundene Entscheidungsgrundlagen
darstellten, an welchen Übergängen eine explizite menschliche Autorisierung
erfolgte und wie nachgelagerte Modellierungs- und
Transformationsentscheidungen an bereits autorisierte Informationszustände
gebunden wurden.

\subsection{Informationsgrundlage und Kontext}
\label{sec:architecture_information_authority}

Die Interpretation heterogener Engineering-Information profitiert in der
Regel von zusätzlichem Kontext. Retrieval-Augmented-Generation-Ansätze
ermöglichen die Einbindung solcher externer Wissensbestände in die
Verarbeitung eines Sprachmodells \cite{Chis.2025}. Semantische Klarstellungen
und Ontologien können ergänzend Begriffe, Relationen und gemeinsame
Konzeptstrukturen repräsentieren und dadurch die korrekte Interpretation der
vorliegenden Informationen unterstützen
\cite{Abonyi.2024,ISO21838_2_2021}. Zudem kann bei der Verarbeitung
heterogener Informationsquellen zusätzliches Kontextwissen erforderlich
sein, um semantische Unterschiede zwischen Quellinformationen zu
berücksichtigen \cite{Bressan.2000}.

Die Nutzung zusätzlicher Information führt dazu, dass Engineering-Dokumente,
Terminologie, Domänenwissen, Ontologien, Modellierungsregeln und technische
Verarbeitungsanweisungen unterschiedliche Rollen einnehmen. Eine explizite
Trennung dieser Rollen macht nachvollziehbar, welche Information eine
fachliche Aussage über das betrachtete System begründet und welche
Information zu ihrer Interpretation oder technischen Verarbeitung beiträgt.
Aus dem Provenienzgedanken, Herkunft und Ableitungsbeziehungen von
Informationsobjekten nachvollziehbar zu erhalten
\cite{MoreauMissier2013}, wird deshalb für diese Arbeit eine zusätzliche
Anforderung abgeleitet: Neben der Herkunft einer Information wird auch ihre
Rolle innerhalb der Modellableitung explizit bestimmt.

Die entwickelte Architektur unterscheidet hierzu drei Informationsrollen:

\begin{itemize}
    \item \emph{Engineering Sources} werden im Folgenden als
    Engineering-Quellen bezeichnet. Sie bilden die fachliche Ausgangsbasis
    für Aussagen über das zu modellierende System. Aus ihnen abgeleitete und
    an ihren konkreten Quellenbezug gebundene Information bildet die
    quellengebundene Informationsgrundlage für nachgelagerte Interpretations-
    und Autorisierungsentscheidungen. Im technischen Modell des Prototyps wird
    diese quellengebundene Information als \emph{source-grounded Evidence}
    bezeichnet.

    \item \emph{Context Sources} werden als Kontextquellen bezeichnet und
    stellen zusätzliche Information für Interpretation und Einordnung bereit.
    Sie können die Interpretation einer Engineering-Quelle unterstützen,
    bilden jedoch nicht selbst die fachliche Ausgangsbasis für Aussagen über
    das betrachtete System.

    \item Der Verarbeitungs- und Orchestrierungskontext umfasst die technische
    Konfiguration und Steuerung des Verarbeitungsvorgangs. Er beeinflusst die
    Art der Verarbeitung, begründet jedoch ebenfalls keine fachliche Aussage
    über das betrachtete System.
\end{itemize}

Tabelle~\ref{tab:information_authority_roles} fasst Funktion und fachliche
Reichweite dieser Informationsrollen zusammen.

\begin{table}[htbp]
\centering
\caption[Informationsrollen der Informationsverarbeitungsarchitektur]{
Informationsrollen und ihre Funktion innerhalb der
Informationsverarbeitungsarchitektur. Engineering-Quellen bilden die
fachliche Ausgangsbasis, Kontextquellen unterstützen deren Interpretation
und der Verarbeitungs- und Orchestrierungskontext bestimmt die technische
Durchführung.}
\label{tab:information_authority_roles}
\small
\begin{tabular}{p{0.23\textwidth} p{0.34\textwidth} p{0.34\textwidth}}
\hline
\textbf{Informationsrolle} &
\textbf{Funktion} &
\textbf{Fachliche Reichweite} \\
\hline

Engineering-Quelle
(\emph{Engineering Source}) &
Fachliche Informationsbasis über das betrachtete System und Ursprung
quellengebundener Information &
Kann fachliche Aussagen über das betrachtete System begründen und bildet
damit die Informationsgrundlage nachgelagerter Interpretations- und
Autorisierungsentscheidungen \\

Kontextquelle
(\emph{Context Source}) &
Unterstützung von Interpretation, Terminologie, Klassifikation und
semantischer Einordnung &
Erweitert die Interpretationsgrundlage; daraus hervorgehende fachliche
Aussagen bleiben an die zugrunde liegende Engineering-Quelle gebunden \\

Verarbeitungs- und
Orchestrierungskontext &
Steuerung und Konfiguration des technischen Verarbeitungsprozesses &
Bestimmt die technische Durchführung und wird als Bestandteil der
Verarbeitungsprovenienz erfasst \\
\hline
\end{tabular}
\end{table}

Zusätzlicher Kontext kann damit gezielt in die Interpretation einbezogen
werden. Terminologie, Domänenwissen oder Ontologien erweitern die
Interpretationsgrundlage, während die daraus hervorgehende fachliche Aussage
auf die zugrunde liegende Engineering-Quelle und die von ihr abgeleitete
quellengebundene Information zurückführbar bleibt. Semantische
Referenzinformation unterstützt damit die fachliche Einordnung, ohne die
Engineering-Quelle als fachliche Ausgangsbasis zu ersetzen
\cite{Abonyi.2024,Chis.2025}.

Dasselbe Prinzip gilt für den technischen Verarbeitungskontext. Prompts,
Agentenrollen oder Ablaufkonfigurationen beeinflussen, wie eine
Engineering-Quelle verarbeitet und wie ein Ergebnis erzeugt wird. Sie werden
deshalb als Bestandteil der Verarbeitungsprovenienz erfasst und machen die
technischen Bedingungen nachvollziehbar, unter denen ein
Verarbeitungsergebnis entstanden ist.

Die Einstufung einer Quelle als Engineering-Quelle legt fest, dass aus ihrem
Inhalt fachliche Information für die weitere Verarbeitung abgeleitet werden
darf. Dies bedeutet noch nicht, dass jeder daraus interpretierte Inhalt
fachlich autorisiert oder für jeden Projektkontext geeignet ist.
Nachgelagerte Interpretations- und Entscheidungsschritte bewerten die
abgeleitete Information und entscheiden über ihre weitere fachliche
Verwendung. Die Informationsrolle bestimmt damit, welche Herkunft eine
fachliche Aussage begründen darf, während die in
Abschnitt~\ref{sec:architecture_human_authority} beschriebene menschliche
Autorisierung über die Verwendung des konkret geprüften Inhalts entscheidet.

Auf dieser kontrollierten Informationsbasis erfolgt die weitere
quellengebundene Interpretation, die im folgenden Abschnitt betrachtet wird.

\subsection{Interpretation von Engineering-Information}
\label{sec:architecture_grounded_interpretation}

Die Einordnung der Quellen aus
Abschnitt~\ref{sec:architecture_information_authority} legt fest, welche
Informationen innerhalb der Verarbeitung welche fachliche Rolle besitzen.
Darauf aufbauend ist zu bestimmen, wie diese Informationen interpretiert
werden. Insbesondere bei mehrdeutigen oder unvollständig formalisierten
Ausgangsinformationen können mehrere fachlich plausible Lesarten bestehen.

Multi-Agent-Ansätze nutzen mehrere Modellinstanzen, um unterschiedliche
Lösungs- oder Argumentationspfade zu erzeugen und miteinander in Beziehung
zu setzen. Für die jeweils untersuchten Aufgaben wurden dadurch unter anderem
Verbesserungen bei Faktentreue und Schlussfolgerungsfähigkeit beobachtet
\cite{Du2024}. Auch die Übereinstimmung mehrerer unabhängig erzeugter
Ergebnisse kann als zusätzliches Bewertungssignal genutzt werden. Sie ist
jedoch kein Nachweis fachlicher Richtigkeit, da die Aussagekraft weiterhin
von den Fähigkeiten der zugrunde liegenden Modelle und von der Gestaltung
ihrer Interaktion abhängt
\cite{Du2024,TriemDing2024,Sunkaraneni2026}.

Im entwickelten Architekturkonzept wird probabilistische KI-Verarbeitung
gezielt dort eingesetzt, wo heterogene oder semantisch offene
Engineering-Information interpretiert werden muss. Mehrere
KI-Perspektiven dienen dazu, unterschiedliche Interpretationen derselben
fachlichen Informationsgrundlage sichtbar und vergleichbar zu machen.
Abweichungen zwischen ihren Ergebnissen können dadurch als
Interpretationsvarianz derselben Ausgangsinformation betrachtet werden.

Die Architektur trennt hierzu zunächst die Identifikation der gemeinsamen
Informationsgrundlage von ihrer multiperspektivischen Interpretation.
Innerhalb einer registrierten Engineering-Quelle wird relevante Information
identifiziert und über konkrete Quellenbezüge an ihre Herkunft gebunden.
Darauf aufbauend werden fachliche Referenzgegenstände
(\emph{Engineering Subjects}) gebildet. Ein Engineering Subject fasst dabei
die quellengebundene Information zusammen, auf die sich die nachfolgenden
Interpretationen beziehen. Es stellt zu diesem Zeitpunkt noch kein
Modellelement des späteren Zielmodells dar.

Auf jeden dieser fachlichen Referenzgegenstände werden anschließend mehrere
konfigurierte Persona-Perspektiven angewendet. Eine \emph{Persona} bezeichnet
in dieser Arbeit die konfigurierte fachliche Perspektive, aus der eine
Agentenrolle eine definierte Verarbeitungsaufgabe bearbeitet. Die
Agentenrolle legt damit die Aufgabe fest, während die Persona die Perspektive
ihrer Bearbeitung bestimmt. Die im Prototyp eingesetzten
Persona-Definitionen sind in
Anhang~\ref{app:persona_definitions} dokumentiert.

Der Verarbeitungspfad folgt damit der Sequenz

\begin{center}
Engineering-Quelle\\
$\downarrow$\\
quellengebundene Information\\
$\downarrow$\\
fachlicher Referenzgegenstand (\emph{Engineering Subject})\\
$\downarrow$\\
mehrere Persona-Interpretationen.
\end{center}

Indem sich alle Persona-Interpretationen auf denselben fachlichen
Referenzgegenstand beziehen, können ihre Ergebnisse unmittelbar miteinander
verglichen werden. Unterschiede zwischen den Perspektiven bleiben dadurch
sichtbar und können gezielt als Interpretationsvarianz ausgewertet werden.

Auf dieser gemeinsamen Grundlage unterscheidet die Architektur mehrere für
die nachgelagerte Bewertung relevante Zustände:

\begin{itemize}
    \item \emph{Übereinstimmung} (\emph{Consensus}) bezeichnet eine
    inhaltliche Übereinstimmung unterschiedlicher Persona-Perspektiven.

    \item \emph{Varianz} (\emph{Variance}) bezeichnet fachlich relevante
    Abweichungen zwischen diesen Perspektiven.

    \item \emph{Gemeinsame Mehrdeutigkeit} (\emph{Shared Ambiguity})
    bezeichnet den Zustand, in dem mehrere Perspektiven zu ähnlichen
    Interpretationen gelangen und zugleich dieselbe Unsicherheit, fehlende
    Information oder Mehrdeutigkeit erkennen.
\end{itemize}

Diese Zustände stellen unterschiedliche Bewertungsinformationen bereit.
Consensus zeigt, dass mehrere Perspektiven zu einer vergleichbaren
Interpretation gelangt sind, ohne damit deren fachliche Richtigkeit
vorwegzunehmen. Variance macht unterschiedliche fachliche Lesarten oder eine
für eine eindeutige Einordnung unzureichende Informationsbasis sichtbar.
Gemeinsame Mehrdeutigkeit zeigt, dass mehrere Perspektiven dieselbe
unaufgelöste Informationslage erkennen. Die multiperspektivische
Verarbeitung stellt damit zusätzliche Entscheidungsgrundlagen für die
nachgelagerte menschliche Prüfung bereit.

Davon zu unterscheiden ist die \emph{Stabilität} einer einzelnen Persona.
Wiederholte Ausführungen derselben Perspektive zeigen, inwieweit die
Interpretation desselben fachlichen Referenzgegenstands zu vergleichbaren
Ergebnissen führt. Die Stabilität beschreibt damit die Reproduzierbarkeit
einer Perspektive. Ihr fachliches Gewicht ergibt sich nicht aus der Anzahl
technischer LLM-Aufrufe.

Consensus, Variance, gemeinsame Mehrdeutigkeit und Stabilität bilden somit
maschinell erzeugte Bewertungsinformationen. Sie machen stabile,
divergierende oder unaufgelöste Interpretationen sichtbar und bilden die
Grundlage für die anschließende fachliche Entscheidung. Die Architektur
trennt damit die probabilistische Interpretation quellengebundener
Engineering-Information von ihrer fachlichen Autorisierung.

Die fachliche Autorisierung bildet den nächsten Entscheidungsschritt im
Transformationsprozess und wird im folgenden Abschnitt im Rahmen der
\emph{Human Authority} betrachtet.

\subsection{Menschliche Entscheidungen im Transformationsprozess}
\label{sec:architecture_human_authority}

Die in Abschnitt~\ref{sec:architecture_grounded_interpretation}
beschriebenen Interpretationsperspektiven erzeugen
Bewertungsinformationen über einen quellengebundenen fachlichen
Referenzgegenstand (\emph{Engineering Subject}). Die fachliche Autorisierung
der daraus hervorgehenden Information erfolgt in einem gesonderten
menschlichen Entscheidungsschritt. Dabei wird festgelegt, welche
interpretierten Inhalte für die weitere fachliche Verwendung im
Transformationsprozess autorisiert werden.

Human-in-the-Loop-Ansätze umfassen unterschiedliche Formen menschlicher
Einbindung, von Überwachung und Rückmeldung bis zur aktiven Bewertung und
Entscheidung \cite{MosqueiraRey2023,Lazaros2026}. Auch bei der
LLM-gestützten Erzeugung von Systems-Engineering- und SysML-v2-Artefakten ist
die fachliche Prüfung ein eigener Bestandteil der Bewertung generierter
Ergebnisse \cite{Topcu2025,Stein2026}. Das entwickelte Architekturkonzept
bildet menschliche Entscheidungen deshalb als explizite
Entscheidungsbefugnisse an definierten Übergängen des
Transformationsprozesses ab.

Der Begriff \emph{Human Authority} bezeichnet in dieser Arbeit die explizite
menschliche Entscheidungsbefugnis an einem definierten fachlichen oder
modellbezogenen Übergang. Eine solche Entscheidung legt fest, welcher
Informations- oder Modellzustand für nachgelagerte Engineering-Schritte
autorisiert wird. \emph{Human Authority} bildet damit den übergeordneten
Begriff für die menschlichen Entscheidungsfunktionen innerhalb des
Transformationsprozesses.

Die erste Autorisierungsgrenze liegt zwischen maschinell erzeugter
Bewertungsinformation und autorisierter Engineering-Information.
Quellengebundene Information, Persona-Interpretationen sowie Informationen
über Consensus, Variance und gemeinsame Mehrdeutigkeit bilden hierfür die
Entscheidungsgrundlage. Eine explizite menschliche Engineering-Entscheidung
legt fest, welche Interpretation beziehungsweise welche daraus abgeleitete
fachliche Aussage zu einem Engineering Subject für die weitere Verwendung
autorisiert wird. Diese Entscheidungsbefugnis wird im Folgenden als
\emph{Human Engineering Authority} bezeichnet.

Die Autorisierung bezieht sich dabei auf den tatsächlich geprüften
Informationszustand des jeweiligen Engineering Subject. Entscheidungen der
Human Engineering Authority referenzieren deshalb einen eindeutig
identifizierten und unveränderlich bestimmbaren Prüfzustand. Eine materielle
Änderung des entscheidungsrelevanten Inhalts führt zu einem neuen
Prüfzustand und einer darauf bezogenen erneuten Autorisierung. Das
Engineering Subject bleibt als fachlicher Referenzgegenstand bestehen,
während die unterschiedlichen Interpretations- und Informationszustände mit
ihrer jeweiligen Provenienz weiterhin referenzierbar bleiben.

Das Human-Authority-Konzept umfasst mehrere menschliche
Entscheidungsgrenzen mit unterschiedlichen Entscheidungsinhalten. Die
\emph{Human Engineering Authority} autorisiert die fachliche Verwendung von
Engineering-Information zu einem Engineering Subject. Weitere Entscheidungen
betreffen die strukturelle Einordnung und Repräsentation dieser Information
im Zielmodell sowie die abschließende Prüfung und Freigabe des erzeugten
Modellzustands.

Für die strukturelle Einordnung im Zielmodell können maschinell erzeugte oder
multiperspektivische Vorschläge als Entscheidungsgrundlage dienen. Die
explizite menschliche Entscheidung legt anschließend fest, in welcher
Modellform und an welcher Stelle des vorgesehenen Modellrahmens die
autorisierte Engineering-Information repräsentiert wird. Erst an dieser
Stelle wird damit festgelegt, durch welches Modellelement beziehungsweise
welche Modellbeziehung der zuvor fachlich autorisierte Inhalt im Zielmodell
abgebildet wird. Diese Entscheidungsbefugnis wird als
\emph{Model Authority} bezeichnet.

Human Engineering Authority und Model Authority adressieren damit
unterschiedliche Entscheidungsgegenstände. Die Human Engineering Authority
entscheidet über die fachliche Verwendung einer interpretierten Information
zu einem Engineering Subject. Die Model Authority entscheidet anschließend
über deren strukturelle Repräsentation im Zielmodell. Die nachgelagerte
Modellbildung verarbeitet damit einen bereits fachlich und strukturell
autorisierten Zustand.

Eine weitere Entscheidungsgrenze liegt am Ende der Modellerzeugung.
Validierung, menschliche Modellprüfung und Freigabe adressieren dabei
aufeinander aufbauende Prüfgegenstände. Die Validierung bewertet den konkret
erzeugten Artefaktzustand anhand der vorgesehenen Invarianten und seiner
Verarbeitbarkeit in der SysML-v2-Zielumgebung. Die anschließende menschliche
Modellprüfung vergleicht den erzeugten Modellzustand mit den zuvor
autorisierten Engineering- und Modellentscheidungen. Die
Freigabeentscheidung autorisiert schließlich genau diesen geprüften und
validierten Zustand zur Veröffentlichung.

Menschliche Entscheidungsbefugnis wird damit gezielt an Übergängen
eingesetzt, die die fachliche Bedeutung, die weitere Verwendung oder die
Modellrepräsentation materiell beeinflussen. Regelgebundene Transformationen
können anschließend auf Basis bereits autorisierter Voraussetzungen und
eindeutig bestimmter Transformationsregeln ausgeführt werden.
Human-in-the-Loop wird damit im vorliegenden Architekturkonzept als gezielte
Zuweisung menschlicher Entscheidungsbefugnis innerhalb des
Transformationsprozesses ausgestaltet
\cite{MosqueiraRey2023,Lazaros2026}.

Der architektonische Zweck der Human Authority besteht darin, explizit und
nachvollziehbar festzuhalten, auf welcher Informationsbasis ein fachlich
oder modellbezogen relevanter Zustandsübergang autorisiert wurde und welche
menschliche Entscheidungsbefugnis dafür zuständig war. Probabilistische
Interpretationen und deren Bewertungsinformationen bilden hierfür
Entscheidungsgrundlagen, während die nachgelagerte Validierung und
Modellprüfung den konkret erzeugten Artefaktzustand adressieren.

Auf die fachliche Autorisierung folgt die Entscheidung, wie die autorisierte
Engineering-Information im Zielmodell strukturell eingeordnet und
repräsentiert wird. Fachliche Bedeutung, strukturelle Einordnung und
konkrete Zielmodellformulierung bilden dabei unterschiedliche
Entscheidungsinhalte, die im folgenden Abschnitt betrachtet werden.

\subsection{Von fachlicher Bedeutung zur Modellrepräsentation}
\label{sec:architecture_meaning_representation}

Mit der Human Engineering Authority aus
Abschnitt~\ref{sec:architecture_human_authority} liegt ein fachlich
autorisierter Informationszustand vor. Er legt fest, welche
Engineering-Aussage zu einem Engineering Subject für die weitere fachliche
Verwendung autorisiert ist. Damit ist jedoch noch nicht entschieden, wie
diese Aussage innerhalb des Systemmodells strukturell eingeordnet und
repräsentiert wird.

Diese Unterscheidung ist insbesondere für semantische und ontologische
Vorgaben relevant. Wie in
Abschnitt~\ref{sec:ki_kontrollmechanismen} beschrieben, stellen Ontologien
eine formalisierte Referenzbasis für Begriffe und Relationen einer Domäne
bereit und unterstützen damit die semantische Integration von
Engineering-Information \cite{OrellanaMandrick2019,Lu2022}.
Domänenspezifische und ontologiebasierte MBSE-Ansätze verdeutlichen darüber
hinaus die Bedeutung gemeinsamer semantischer Bezugssysteme für die
strukturierte Nutzung von Engineering-Information
\cite{Liu2025SysDO,Dong2026}.

Für die Zielmodellbildung wird diese semantische Einordnung um
modellbezogene Entscheidungen ergänzt. Der Modellierungsrahmen, seine
strukturellen Konventionen, die vorgesehenen Sichten und der jeweilige
Modellierungszweck bestimmen, welche Repräsentationen für eine fachlich
autorisierte Information grundsätzlich in Betracht kommen. Die fachliche
Bedeutung einer Information und ihre konkrete Repräsentation im Zielmodell
werden dadurch als getrennte Entscheidungsinhalte behandelt.

Die Architektur des in dieser Arbeit vorgestellten Transformationssystems
unterscheidet hierzu drei aufeinander aufbauende Entscheidungs- und
Verarbeitungsebenen:

\begin{itemize}
    \item Die \emph{fachliche Bedeutung} (\emph{Engineering Meaning})
    beschreibt die durch die autorisierte Engineering-Information getragene
    Aussage.

    \item Die \emph{Einordnung und Repräsentation im Zielmodell} legt fest,
    durch welches Modellelement beziehungsweise welche Modellbeziehung und an
    welcher Stelle des vorgesehenen Modellrahmens diese Information
    repräsentiert wird.

    \item Die \emph{zielmodellspezifische Formulierung} bestimmt, wie der
    autorisierte Modellinhalt innerhalb der unterstützten Zielnotation
    ausgedrückt wird.
\end{itemize}

Abbildung~\ref{fig:architecture_target_model_formation} stellt den daraus
resultierenden Übergang von fachlich autorisierter Engineering-Information
zur SysML-v2-Repräsentation zusammenfassend dar.

\begin{wrapfigure}{r}{0.47\textwidth}
    \centering
    \vspace{-0.8\baselineskip}
    \includegraphics[
        width=\linewidth,
        keepaspectratio
    ]{Graphiken/4.5_Zielmodellbildung_halbseite_v5.pdf}
    \caption[Kontrollierte Zielmodellbildung]{
        Kontrollierte Zielmodellbildung von fachlich autorisierter
        Engineering-Information bis zur SysML-v2-Repräsentation.
        Semantische und ontologische Vorgaben, Modellierungsrahmen und
        strukturelle Konventionen sowie Vorgaben der Zielnotation unterstützen
        die jeweils zugehörigen Verarbeitungsschritte.
        Eigene Darstellung.}
    \label{fig:architecture_target_model_formation}
    \vspace{-0.6\baselineskip}
\end{wrapfigure}

Aus der fachlich autorisierten Engineering-Information werden zunächst
mögliche Modellrepräsentationen abgeleitet. Eine fachliche Aussage kann dabei
beispielsweise als Anforderung, Funktion, logische Komponente oder Beziehung
im Zielmodell repräsentiert werden. Welche dieser Repräsentationen geeignet
ist und an welcher Stelle des projektspezifischen Modellrahmens sie
eingeordnet wird, ist Gegenstand einer eigenen modellbezogenen Entscheidung.
Die SysML-v2-Spezifikation definiert die hierfür verfügbaren sprachlichen
Konzepte und deren formale Bedeutung \cite{SysMLv2}. Die Auswahl eines für
die konkrete Engineering-Aussage geeigneten Konstrukts erfolgt innerhalb des
projektspezifischen Modellierungsprozesses.

Für diese modellbezogene Entscheidung werden aus autorisierter
Engineering-Information \emph{Model Candidates} abgeleitet. Sie beschreiben
mögliche Modellelemente, Beziehungen und strukturelle Zuordnungen, ohne
bereits selbst den autorisierten Modellzustand darzustellen. Der Übergang von
fachlicher Bedeutung zu Modellstruktur wird dadurch als eigener prüfbarer
Verarbeitungsschritt abgebildet.

Die Model Candidates bilden die Entscheidungsgrundlage für die
\emph{Model Authority}. Unterschiedliche Modellierungsperspektiven können
hierzu Vorschläge für die strukturelle Einordnung und Repräsentation
bereitstellen. Die explizite menschliche Entscheidung legt anschließend fest,
welches Modellelement beziehungsweise welche Modellbeziehung verwendet wird
und an welcher Stelle des Modellrahmens die autorisierte
Engineering-Information repräsentiert wird.

Human Engineering Authority und Model Authority adressieren damit
unterschiedliche Entscheidungsinhalte. Die Human Engineering Authority
autorisiert die fachliche Bedeutung und weitere Verwendung der Information zu
einem Engineering Subject. Die Model Authority autorisiert anschließend ihre
strukturelle Einordnung und Repräsentation im Zielmodell.

Nach dieser modellbezogenen Autorisierung werden die autorisierten
Modellinhalte zu einem gemeinsamen \emph{Internal Engineering Model} (IEM)
zusammengeführt. Das IEM bildet, wie in
Abbildung~\ref{fig:architecture_target_model_formation} dargestellt, einen
kontrollierten Modellzustand zwischen autorisierter Engineering-Information
und konkreter Zielnotation. Es bündelt die autorisierten Modellinhalte unter
Erhalt ihrer fachlichen Bedeutung und stellt sie auf einer internen
Repräsentationsebene für die anschließende zielmodellspezifische Formulierung
bereit.

Auf Basis dieses internen Modellzustands wird anschließend festgelegt, wie
Modellelemente und Beziehungen innerhalb des unterstützten Ausschnitts der
SysML-v2-Zielnotation formuliert werden. Fachliche Bedeutung,
Modellorganisation und die Erzeugung der SysML-v2-Repräsentation bleiben
dadurch als explizite und nachvollziehbare Verarbeitungsebenen erhalten.

Für die dritte Teilforschungsfrage ist diese Struktur besonders relevant.
Ontologische und semantische Vorgaben strukturieren Bedeutungsräume,
Terminologie und Relationen
\cite{OrellanaMandrick2019,Lu2022,Dong2026}. Die Human Engineering Authority
autorisiert die fachliche Bedeutung und Verwendung der Engineering-
Information, die Model Authority deren strukturelle Repräsentation und die
zielmodellspezifische Formulierung überführt den autorisierten Modellzustand
in die unterstützte SysML-v2-Notation. Die expliziten Übergänge zwischen
diesen Ebenen bilden damit die Leitplanken für die kontrollierte
Zielmodellbildung.

Der folgende Abschnitt betrachtet darauf aufbauend die regelgebundene und
reproduzierbare Ausführung der nachgelagerten Transformation.

\subsection{Regelgebundene Transformation}
\label{sec:architecture_controlled_transformation}

Mit den in
Abschnitt~\ref{sec:architecture_human_authority} und
Abschnitt~\ref{sec:architecture_meaning_representation}
beschriebenen Autorisierungsentscheidungen liegt ein Zustand vor, dessen
fachliche Bedeutung sowie strukturelle Einordnung und Repräsentation
autorisiert sind. Die Human Engineering Authority autorisiert dabei die
fachliche Verwendung der Engineering-Information, während die Model
Authority deren Repräsentation im Zielmodell festlegt. Die nachgelagerte
Transformation verarbeitet damit bereits fachlich und modellbezogen
autorisierte Inhalte. Ihre Aufgabe besteht darin, diese regelgebunden und
reproduzierbar in die unterstützte SysML-v2-Zielrepräsentation zu
überführen.

Maschinenlesbare und versionskontrollierte Architekturartefakte bilden eine
geeignete Grundlage für automatisierte Transformationen.
\emph{Architecture as Code}- und \emph{Everything as Code}-Ansätze verbinden
maschinenlesbare Repräsentationen unter anderem mit Versionierung,
Automatisierung und reproduzierbarer Weiterverarbeitung
\cite{Bucaioni2025,Stirbu2022}. Das entwickelte Transformationssystem nutzt
diese Eigenschaften, indem nach erfolgter Autorisierung explizite und
versionierte Transformationsregeln auf den jeweiligen Modellzustand
angewendet werden.

Die regelgebundene Verarbeitung beginnt mit der Bildung des
\emph{Internal Engineering Model} (IEM). Autorisierte Modellinhalte werden
hierzu zu einem gemeinsamen internen Modellzustand zusammengeführt. Das IEM
bewahrt die zuvor autorisierte fachliche Bedeutung und strukturelle
Einordnung und stellt sie in einer einheitlichen internen Repräsentation für
die anschließende Zielabbildung bereit. Die in
Abbildung~\ref{fig:architecture_target_model_formation} dargestellte
zielmodellspezifische Formulierung überführt diesen Modellzustand
anschließend in die unterstützte SysML-v2-Repräsentation.

Die hierfür erforderlichen Regeln adressieren unterschiedliche Aspekte der
Transformation. Ein \emph{Target Notation Profile} bestimmt den unterstützten
Ausschnitt der Zielnotation. Ein \emph{Generation Profile} definiert die
Zuordnung der internen Modellkonzepte zu den vorgesehenen Konstrukten der
Zielnotation. Ein \emph{Artifact Structure Profile} legt die Organisation
des erzeugten Artefakts fest, beispielsweise hinsichtlich Packages, Dateien
oder Reihenfolge. Die verwendeten Profile werden versioniert und mit dem
verarbeiteten Modellzustand verknüpft. Dadurch bleibt nachvollziehbar, unter
welchen Transformationsbedingungen ein konkretes SysML-v2-Artefakt erzeugt
wurde.

Die regelgebundene Transformation setzt voraus, dass die erforderlichen
Autorisierungsentscheidungen vorliegen und für den jeweiligen Übergang eine
eindeutig anwendbare Transformationsregel verfügbar ist. Für Fälle, in denen
diese Voraussetzungen nicht vollständig erfüllt sind, verwendet das
Transformationssystem das Prinzip des \emph{fail-closed processing}. Darunter
wird in dieser Arbeit eine Verarbeitung verstanden, bei der ein
Zustandsübergang nur bei vollständig erfüllten Voraussetzungen ausgeführt
wird. Andernfalls entsteht ein expliziter blockierender Zustand, der die
fehlende Voraussetzung sichtbar macht und eine weitere Verarbeitung
verhindert.

Ein solcher blockierender Zustand kann beispielsweise durch eine fehlende
Autorisierung, widersprüchliche Identitäten oder Provenienzbindungen,
nicht aufgelöste technische Referenzen oder eine fehlende beziehungsweise
nicht eindeutig anwendbare Zielabbildung ausgelöst werden. Für einen nicht
unterstützten Beziehungstyp wird beispielsweise keine
SysML-v2-Repräsentation erzeugt, solange keine dafür definierte und
anwendbare Abbildungsregel vorliegt. Ist stattdessen eine zusätzliche
fachliche oder modellbezogene Entscheidung erforderlich, wird die
Verarbeitung an der dafür vorgesehenen Autorisierungsgrenze fortgesetzt.

Dasselbe Prinzip gilt für die Bildung des internen Modellzustands. Die
Modellassemblierung erfolgt auf Basis bereits autorisierter Modellinhalte und
definierter struktureller Regeln. Erfordert die Integration eines
Modellinhalts eine zusätzliche fachliche oder modellbezogene Entscheidung,
wird der Verarbeitungspfad an dieser Stelle unterbrochen. Die offene
Entscheidung wird der zuständigen Human Authority zugeführt und die
regelgebundene Verarbeitung erst auf Grundlage eines entsprechend
autorisierten Zustands fortgesetzt.

Das Transformationssystem priorisiert damit einen kontrollierten und
nachvollziehbaren Verarbeitungszustand gegenüber einer vollständigen
Modellerzeugung unter unklaren Voraussetzungen. Ein SysML-v2-Artefakt wird
nur für Inhalte erzeugt, für die die erforderlichen Autorisierungen und
Transformationsregeln vorliegen. Ein blockierender Zustand ist deshalb ein
vorgesehener Zustand des Transformationsprozesses und kennzeichnet eine noch
ausstehende Entscheidung oder Transformationsvoraussetzung.

An die Generierung schließt sich die Validierung des konkret erzeugten
Artefaktzustands an. Auch Arbeiten zur LLM-gestützten
SysML-v2-Erzeugung unterscheiden zwischen der Erzeugung eines Artefakts und
seiner nachgelagerten Prüfung \cite{Stein2026}. Die Validierung bewertet den
erzeugten Zustand anhand definierter Invarianten und seiner vorgesehenen
Verarbeitbarkeit in der SysML-v2-Zielumgebung. Das Ergebnis der Validierung
wird an genau den geprüften Artefaktzustand gebunden.

Ein erfolgreich validierter Artefaktzustand kann anschließend in die
nachgelagerte menschliche Modellprüfung und Freigabe übergehen.
Festgestellte Validierungsabweichungen oder eine nicht vollständig
ausführbare Validierung führen dagegen zu einem blockierenden Zustand.
Erforderliche Änderungen erzeugen einen neuen Modell- beziehungsweise
Artefaktzustand, der anschließend erneut validiert wird. Generierung,
Validierung und eine gegebenenfalls erforderliche Überarbeitung bleiben
dadurch als getrennte und nachvollziehbare Verarbeitungsschritte erhalten.

Das entwickelte Transformationssystem verbindet damit probabilistische und
regelgebundene Verarbeitung an unterschiedlichen Stellen des
Transformationsprozesses. Probabilistische Verfahren unterstützen die
Interpretation, den Vergleich und die Erzeugung von
Entscheidungsalternativen. Fachlich oder modellbezogen relevante Ergebnisse
werden durch die jeweils zuständige Human Authority autorisiert. Liegen die
erforderlichen Autorisierungen und eindeutigen Transformationsregeln vor,
erfolgt die weitere Verarbeitung regelgebunden. Fehlende Voraussetzungen
führen zu einem expliziten Haltepunkt und werden, soweit eine menschliche
Entscheidung erforderlich ist, an die zuständige Autorisierungsgrenze
zurückgeführt.

Damit autorisierte Engineering-Information bis zum validierten
Modellartefakt nachvollziehbar weiterverarbeitet werden kann, müssen die
Beziehungen zwischen Quellen, Interpretationen,
Autorisierungsentscheidungen, Modellzuständen und Zielartefakten über den
gesamten Transformationsprozess erhalten bleiben. Diese Querschnittsfunktion
wird im folgenden Abschnitt durch Provenienz und Nachverfolgbarkeit
betrachtet.

\subsection{Durchgängige Provenienz und Nachverfolgbarkeit}
\label{sec:architecture_traceability}

Die bisher beschriebenen Architekturmechanismen adressieren unterschiedliche
Aspekte des Transformationsprozesses. Damit ihre Auswirkungen über den
gesamten Verarbeitungspfad nachvollziehbar bleiben, müssen auch die
Beziehungen zwischen den entstehenden Informations-, Entscheidungs- und
Modellzuständen erhalten werden.

Provenienz beschreibt grundsätzlich die Herkunft und
Entstehungszusammenhänge von Information. Das W3C PROV Data Model modelliert
hierzu unter anderem Entitäten, Aktivitäten und Agenten sowie deren
Ableitungsbeziehungen \cite{MoreauMissier2013}. Provenienzansätze für
datenverarbeitende Prozesse adressieren entsprechend die nachvollziehbare
Herkunft und Transformation von Daten über mehrere Verarbeitungsschritte
hinweg \cite{Simmhan2005}. Im Systems Engineering ergänzt
Nachverfolgbarkeit diese Perspektive durch explizite Beziehungen zwischen
Engineering-Artefakten und deren Entwicklung. Dies ist insbesondere im
Requirements Engineering etabliert \cite{ISO29148}. Auch im
\emph{Architecture as Code}-Kontext werden Versionierung und
Nachverfolgbarkeit als wesentliche Eigenschaften maschinenlesbarer
Architekturartefakte hervorgehoben \cite{Bucaioni2025}.

Für den in dieser Arbeit betrachteten Transformationsprozess umfasst
Nachverfolgbarkeit mehr als die direkte Beziehung zwischen einer
Engineering-Quelle und der resultierenden SysML-v2-Repräsentation. Zwischen
beiden liegen quellengebundene Informationszustände, maschinelle
Interpretationen, menschliche Autorisierungsentscheidungen,
Modellentscheidungen und regelgebundene Transformationen. Die Beziehungen
zwischen diesen Zuständen bilden deshalb einen wesentlichen Bestandteil der
Informationsverarbeitungsarchitektur.

Provenienz und Nachverfolgbarkeit werden hierzu als Querschnittsfunktion des
Transformationssystems behandelt. Jeder für die Modellbildung relevante
Zustandsübergang bleibt mit den Informations- und Entscheidungszuständen
verbunden, auf deren Grundlage er entstanden ist. Dies betrifft sowohl
automatisierte Verarbeitungsschritte als auch menschliche
Autorisierungsentscheidungen.

Abbildung~\ref{fig:architecture_traceability} veranschaulicht diese
Querschnittsfunktion entlang der in
Abschnitt~\ref{sec:architecture_meaning_representation} eingeführten
Zielmodellbildung. Den wesentlichen Informations-, Entscheidungs- und
Modellzuständen sind jeweils die Provenienz- und
Nachverfolgbarkeitsinformationen zugeordnet, die ihre Herkunft, Ableitung und
weitere Verarbeitung rekonstruierbar machen.

\begin{figure}[htbp]
    \centering
    \includegraphics[
        width=0.78\textwidth,
        keepaspectratio
    ]{Graphiken/4.7_Provenienz_und_Nachverfolgbarkeit_halbseite_v4.pdf}
    \caption[Durchgängige Provenienz und Nachverfolgbarkeit]{
        Durchgängige Provenienz- und Nachverfolgbarkeitsbeziehungen entlang
        der Zielmodellbildung. Eigene Darstellung.
    }
    \label{fig:architecture_traceability}
\end{figure}

Die in Abbildung~\ref{fig:architecture_traceability} dargestellten
Zuordnungen erhalten die Beziehungen zwischen den einzelnen Zuständen. Ein
erzeugtes Modellelement oder eine erzeugte Modellbeziehung kann dadurch
sowohl auf die zugrunde liegende fachliche Ausgangsinformation als auch auf
die für ihre weitere Verarbeitung relevanten Autorisierungs- und
Modellentscheidungen zurückgeführt werden. Die verschiedenen Informations-,
Entscheidungs- und Modellzustände bleiben dabei als eigenständige
referenzierbare Zustände erhalten und werden über explizite Beziehungen
miteinander verbunden.

Hierfür verwendet das Transformationssystem stabile Identitäten, explizite
Vorgängerreferenzen und an konkrete Informationszustände gebundene
Entscheidungsreferenzen. Eine menschliche Autorisierung erzeugt einen eigenen
referenzierbaren Entscheidungszustand, der mit dem tatsächlich geprüften
Informations- oder Modellzustand verbunden bleibt. Dieselbe Logik setzt sich
über Modellentscheidungen und interne Modellzustände bis zur
Zielrepräsentation fort.

Die in Abbildung~\ref{fig:architecture_traceability} dargestellten
Beziehungen können in zwei komplementären Richtungen verfolgt werden. Die
Vorwärtsrichtung zeigt, welche nachgelagerten Informations-,
Entscheidungs- und Modellzustände aus einer gegebenen Engineering-Quelle
beziehungsweise einem quellengebundenen Informationszustand hervorgegangen
sind. Die Rückwärtsrichtung zeigt, auf welcher Informations- und
Entscheidungsgrundlage ein konkretes Modellelement oder eine konkrete
Modellbeziehung des resultierenden Modells beruht.

Diese Beziehungen bleiben auch erhalten, wenn mehrere Informationszustände
zu einem gemeinsamen Modellelement oder einer gemeinsamen Modellbeziehung
beitragen. Ein zusammengeführter Zustand referenziert die fachlich
beitragenden Vorgängerzustände und erhält deren jeweilige Herkunft. Dadurch
können Informationen gemeinsam weiterverarbeitet werden, während ihre
quellenbezogenen Provenienzbeziehungen bestehen bleiben.

Die Nachverfolgbarkeit wird bis zur SysML-v2-Zielrepräsentation fortgeführt.
Das \emph{Internal Engineering Model} erhält die Beziehungen seiner Elemente
und Verbindungen zu den autorisierten Engineering- und Modellzuständen. Bei
der regelgebundenen SysML-v2-Generierung werden diese Beziehungen mit den
erzeugten Zielrepräsentationen verknüpft. Dadurch bleibt auch die Erzeugung
der textuellen SysML-v2-Repräsentation als eigener Transformationsschritt
nachvollziehbar.

Auch die Validierung wird an den konkret geprüften Artefaktzustand gebunden.
Ein Validierungsergebnis bezieht sich damit auf einen eindeutig
identifizierten generierten Modellstand. Festgestellte Abweichungen können
über die bestehenden Provenienz- und Nachverfolgbarkeitsbeziehungen den
zugehörigen Modell-, Autorisierungs- und Informationszuständen zugeordnet
werden.

Die anschließende menschliche Modellprüfung und die davon getrennte
Freigabeentscheidung beziehen sich ebenfalls auf den jeweils konkret
geprüften Modellzustand. Wird dieser Zustand verändert, entsteht ein neuer
prüf- und freigabebedürftiger Zustand. Dadurch bleibt der freigegebene
Modellstand eindeutig mit seiner Validierung, der menschlichen Modellprüfung
und den zugehörigen Autorisierungsentscheidungen verbunden.

Das Ergebnis des Transformationsprozesses umfasst damit neben der
SysML-v2-Repräsentation auch eine rekonstruierbare Herleitung des
Modellzustands. Nachvollziehbar bleiben die zugrunde liegenden
Engineering-Quellen und quellengebundenen Informationszustände, maschinellen
Interpretationen, menschlichen Autorisierungsentscheidungen,
Modellentscheidungen sowie Validierungs- und Freigabezustände.

Der zentrale Architekturanspruch des entwickelten Transformationssystems
besteht darin, die fachliche und modellbezogene Herleitung eines
resultierenden Modellzustands über den gesamten Transformationsprozess
nachvollziehbar zu erhalten. Provenienz und Nachverfolgbarkeit verbinden
hierzu die zuvor getrennten Informations-, Interpretations-,
Autorisierungs-, Modellierungs- und Transformationsschritte zu einem
durchgängigen Informationsfluss.

Diese Querschnittsfunktion gewinnt zusätzliche Bedeutung, sobald mehrere
Engineering-Quellen gemeinsam zu einem Modell beitragen. Dann müssen
Informationen projektweit zusammengeführt werden können, während ihre
quellenbezogene Herkunft und die daran gebundenen Entscheidungen erhalten
bleiben. Die daraus entstehenden Anforderungen werden im folgenden Abschnitt
betrachtet.

\subsection{Verarbeitung mehrerer Quellen}
\label{sec:architecture_multisource}

Die bisher entwickelten Architekturprinzipien wurden zunächst entlang der
Verarbeitung einzelner Engineering-Quellen hergeleitet. In realen
Engineering-Projekten verteilt sich relevante Information jedoch regelmäßig
auf mehrere Dokumente, Datenbestände und fachliche Sichten. Für die
Modellbildung entsteht daraus die Anforderung, Informationen aus mehreren
Quellen gemeinsam zu berücksichtigen und gleichzeitig ihre jeweilige
Herkunft und Autorisierung nachvollziehbar zu erhalten.

Semantische Modelle und Ontologien können die Vergleichbarkeit solcher
Informationen unterstützen, indem sie gemeinsame Begriffe, Relationen und
Konzeptstrukturen bereitstellen
\cite{OrellanaMandrick2019,Lu2022,Dong2026}. LLM-gestützte semantische
Zuordnungen können durch formale Taxonomien beziehungsweise ontologische
Referenzen strukturiert werden, sind jedoch gegenüber Fehlzuordnungen und
Überspezifikation anfällig \cite{Souza.2026}. Semantische Ähnlichkeit bildet
damit einen Hinweis auf mögliche fachliche Zusammenhänge, deren Bedeutung
im jeweiligen Projektkontext bewertet werden muss.

Für die Verarbeitung mehrerer Quellen ergeben sich zwei miteinander
verbundene Anforderungen. Die Provenienz jeder Information bleibt über den
gesamten Transformationsprozess erhalten, während fachlich zusammengehörige
Informationen projektweit gemeinsam betrachtet werden können. Dadurch können
Informationen unterschiedlicher Herkunft zusammengeführt werden, ohne ihre
quellenbezogenen Beziehungen und Autorisierungszustände aufzulösen.

Das Transformationssystem verarbeitet Informationen zunächst
quellenbezogen. Registrierung, quellengebundene Informationszustände,
Interpretation und die daraus hervorgehenden autorisierten
Engineering-Informationen bleiben mit der jeweiligen Engineering-Quelle
verbunden. Diese quellenbezogene Verarbeitung schafft die Grundlage dafür,
Beiträge mehrerer Quellen anschließend projektweit miteinander in Beziehung
zu setzen.

Vor dieser projektweiten Zusammenführung wird geprüft, ob eine Quelle für
den betrachteten System- und Projektkontext relevant und für die
projektweite Verwendung geeignet ist. Diese
\emph{projektbezogene Quellenzulassung} (\emph{Project Fit}) bildet eine
eigene Kontrollgrenze des Transformationsprozesses. Für ihre Bewertung
können vorhandene Projektinformationen, Terminologie, Kontextquellen und
bereits autorisierte Engineering-Informationen herangezogen werden.

Das Ergebnis der Project-Fit-Bewertung dient als Entscheidungsgrundlage für
die projektweite Verwendung einer Quelle. Eindeutig zuordenbare Quellen
können entsprechend der festgelegten Projektregeln weiterverarbeitet werden.
Bei unklarer oder widersprüchlicher Zuordnung wird eine explizite
menschliche Entscheidung über ihre Verwendung im Projektkontext
erforderlich. Die Einordnung einer Quelle als Engineering-Quelle oder
Kontextquelle bleibt davon getrennt und wird weiterhin als eigenständige
Informationsrolle behandelt.

Nach der Quellenzulassung werden Beziehungen zwischen autorisierten
Engineering-Informationen unterschiedlicher Quellen in einem
\emph{projektweiten semantischen Abgleich} untersucht. Das
Transformationssystem unterscheidet hierfür folgende Beziehungsarten:

\begin{itemize}
    \item \emph{equivalent}: mehrere Beiträge drücken im Wesentlichen
          dieselbe fachliche Bedeutung aus,

    \item \emph{complementary}: Beiträge behandeln denselben fachlichen
          Zusammenhang und ergänzen sich,

    \item \emph{potential conflict}: Beiträge enthalten möglicherweise
          unvereinbare Aussagen zum selben fachlichen Gegenstand,

    \item \emph{distinct}: Beiträge behandeln unterschiedliche fachliche
          Gegenstände,

    \item \emph{uncertain}: die Beziehung kann nicht ausreichend eindeutig
          bestimmt werden.
\end{itemize}

Diese Kategorien strukturieren die projektweite Gegenüberstellung
autorisierter Engineering-Informationen. Ontologische und terminologische
Referenzen unterstützen dabei den semantischen Vergleich
\cite{OrellanaMandrick2019,Lu2022}. Die daraus abgeleiteten Beziehungen
bilden eine Entscheidungsgrundlage für die weitere fachliche
Zusammenführung.

Erfordert die Zusammenführung eine neue oder materiell veränderte
Engineering-Aussage, erfolgt eine erneute fachliche Autorisierung. Die
\emph{Human Engineering Authority} entscheidet dabei auf projektweiter
Ebene, ob Informationen unabhängig fortbestehen, gemeinsam verwendet
werden, einander ablösen oder zunächst als ungeklärte Beziehung erhalten
bleiben. Die quellenbezogen autorisierten Ausgangszustände und ihre
Provenienz bleiben parallel dazu referenzierbar.

Damit wird die Human Engineering Authority auch bei der Verarbeitung
mehrerer Quellen für dieselbe Art fachlicher Entscheidung eingesetzt wie
zuvor: Sie autorisiert die Verwendung von Engineering-Information. Der
Unterschied liegt in der Informationsbasis der Entscheidung, die nun
Beiträge aus mehreren bereits quellenbezogen autorisierten Zuständen
umfasst. Eine zusätzliche Kategorie menschlicher Entscheidungsbefugnis ist
hierfür nicht erforderlich.

Projektweit autorisierte Engineering-Information kann anschließend
Auswirkungen auf einen bestehenden Modellzustand besitzen. Zunächst wird
fachlich entschieden, ob und wie sich der bisher autorisierte
Engineering-Zustand verändert. Darauf aufbauend untersucht ein getrennter
\emph{Modellwirkungsabgleich}, welche Auswirkungen die fachliche Änderung auf
den bestehenden Modellzustand hat.

Damit bleibt die in
Abschnitt~\ref{sec:architecture_meaning_representation} entwickelte Trennung
zwischen Human Engineering Authority und Model Authority auch bei der
Verarbeitung mehrerer Quellen erhalten. Eine neue Engineering-Aussage kann
eine bestehende Modellrepräsentation bestätigen, ergänzen, verändern oder
ersetzen. Die \emph{Model Authority} entscheidet auf Grundlage der zuvor
autorisierten fachlichen Änderung über die daraus resultierende strukturelle
Einordnung und Repräsentation im Zielmodell.

Die Verarbeitung mehrerer Quellen bildet damit einen Anwendungsfall der
zuvor entwickelten Architekturprinzipien auf projektweiter Ebene.
Quellenidentität, Provenienz, quellenbezogene Autorisierungszustände, Human
Engineering Authority und Model Authority bleiben auch bei der
Zusammenführung von Informationen unterschiedlicher Herkunft erhalten. Das
Transformationssystem verbindet dadurch quellenbezogene Nachvollziehbarkeit
mit projektweiter fachlicher Zusammenführung.

Diese Kombination bildet die Grundlage für das im folgenden Abschnitt
zusammengeführte Architekturkonzept.

\subsection{Zusammenführung zum Architekturkonzept}
\label{sec:architecture_integrated_concept}

Die vorangegangenen Abschnitte haben die zentralen Mechanismen des
Architekturkonzepts schrittweise hergeleitet. In ihrer Gesamtheit bilden sie
eine Informationsverarbeitungsarchitektur zur nachvollziehbaren Überführung
heterogener Engineering-Information in einen autorisierten und validierbaren
SysML-v2-Modellzustand.

Das Architekturkonzept verbindet probabilistische und regelgebundene
Verarbeitung entsprechend der jeweiligen Aufgabenstellung.
Probabilistische Verfahren unterstützen semantische Interpretation, Vergleich
und die Erzeugung von Entscheidungsalternativen. Fachlich oder modellbezogen
relevante Zustandsübergänge werden durch explizite menschliche
Autorisierungsentscheidungen kontrolliert
\cite{Topcu2025,Stein2026}. Die Human Engineering Authority autorisiert die
fachliche Verwendung von Engineering-Information, während die Model Authority
deren strukturelle Einordnung und Repräsentation im Zielmodell festlegt.
Liegen die erforderlichen Autorisierungen und eindeutige
Transformationsregeln vor, erfolgt die weitere Verarbeitung regelgebunden.
Für unvollständige oder nicht eindeutig erfüllte
Transformationsvoraussetzungen greift das in
Abschnitt~\ref{sec:architecture_controlled_transformation} eingeführte
\emph{fail-closed processing}.

Abbildung~\ref{fig:integrated_architecture} ergänzt diese prozessbezogene
Beschreibung um eine Use-Case-Sicht des integrierten Architekturkonzepts.
Die dargestellten Anwendungsfälle bündeln die zuvor beschriebenen Mechanismen
aus Sicht ihrer Nutzung: Bewertung und Vorbereitung von Engineering-Quellen,
Bewertung semantischer Konsistenz und Behandlung von Mehrdeutigkeit, menschliche
Autorisierung von Zwischenergebnissen, projektbezogene Quellenzulassung,
SysML-v2-Erzeugung und Validierung sowie die abschließende Modellprüfung und
Freigabe. Die innerhalb der Use Cases ausgewiesenen Akteursrollen
kennzeichnen die jeweils beteiligten menschlichen Rollen. Sie sind von den in
Abschnitt~\ref{sec:architecture_human_authority} beschriebenen
Entscheidungsbefugnissen der Human Authority zu unterscheiden. Die
Use-Case-Sicht beschreibt damit die Interaktion mit den
Architekturmechanismen und keine vollständige zeitliche Ablaufreihenfolge.

Die in der Use-Case-Sicht verwendeten Akteursrollen repräsentieren
unterschiedliche menschliche Aufgabenbereiche innerhalb des
Transformationsprozesses. Der \emph{SystemModeler} adressiert insbesondere
die formale Modellbildung und die strukturelle Verarbeitung von
Modellinformation. Der \emph{DomainExpert} bringt die für die fachliche
Interpretation und Klärung mehrdeutiger Information erforderliche
Domänenkenntnis ein. Der \emph{SystemArchitect} vertritt die
modellbezogene und architektonische Perspektive, insbesondere bei der
Prüfung und Einordnung des resultierenden Systemmodells. Der
\emph{Information Manager} adressiert die Qualität, Herkunft und
semantische Konsistenz der verarbeiteten Information. Die Rollen wurden
damit so gewählt, dass die fachliche, informationsbezogene und
modellbezogene menschliche Beteiligung innerhalb des Architekturkonzepts
abgebildet werden kann. Sie sind nicht mit den in
Abschnitt~\ref{sec:architecture_human_authority} definierten
Entscheidungsbefugnissen gleichzusetzen: Eine Rolle kann abhängig vom
jeweiligen Anwendungsfall eine entsprechende Human Authority ausüben.

\begin{figure}[htbp]
    \centering
    \includegraphics[
        width=\textwidth,
        keepaspectratio
    ]{Graphiken/4.9_UseCase.png}
    \caption[Use-Case-Sicht des integrierten Architekturkonzepts]{
        Use-Case-Sicht des integrierten Architekturkonzepts. Die Use Cases
        repräsentieren die wesentlichen Anwendungsfälle von der Bewertung und
        Vorbereitung einer Engineering-Quelle über semantische Bewertung,
        menschliche Autorisierung und projektbezogene Quellenzulassung bis zur
        SysML-v2-Erzeugung, Validierung sowie abschließenden Modellprüfung und
        Freigabe. Die innerhalb der Use Cases angegebenen Akteursrollen zeigen
        die jeweils beteiligten menschlichen Rollen. Beziehungen zum
        Anwendungsfall \emph{ReviewAndAuthorizeIntermediateResult}
        kennzeichnen die explizite Einbindung einer menschlichen
        Zwischenautorisierung. Eigene Darstellung.}
    \label{fig:integrated_architecture}
\end{figure}

Provenienz und Nachverfolgbarkeit wirken als Querschnittsfunktion über den
gesamten Transformationsprozess. Sie erhalten die Beziehungen zwischen
Quellen, maschinellen Verarbeitungsergebnissen, menschlichen
Autorisierungsentscheidungen, Modellzuständen und Zielartefakten. Dadurch
bleibt die Herleitung eines resultierenden Modellzustands über seine
fachlichen und modellbezogenen Vorgängerzustände rekonstruierbar
\cite{MoreauMissier2013}.

Die Verarbeitung mehrerer Engineering-Quellen erweitert das
Architekturkonzept auf eine projektweite Ebene. Quellenbezogene Identität,
Provenienz und Autorisierungszustände bleiben erhalten, während zugelassene
Informationen semantisch miteinander in Beziehung gesetzt werden. Erfordert
ihre Zusammenführung eine neue oder materiell veränderte
Engineering-Aussage, erfolgt eine projektweite Anwendung der Human
Engineering Authority. Die daraus autorisierte Engineering-Information kann
anschließend über die Model Authority in den bestehenden Modellzustand
integriert werden. Frühere Informations-, Entscheidungs- und Modellzustände
bleiben dabei als Teil der Provenienzkette referenzierbar.

Der wissenschaftliche Beitrag des Architekturkonzepts liegt in der
systematischen Verbindung einer kontrollierten Informationsbasis,
quellengebundener multiperspektivischer Interpretation, expliziter
menschlicher Autorisierung, der Trennung von Engineering Meaning und
Zielmodellbildung, regelgebundener und fail-closed kontrollierter
Transformation sowie durchgängiger Provenienz und Nachverfolgbarkeit. Diese
Mechanismen bilden gemeinsam eine konsistente
Informationsverarbeitungsarchitektur für die KI-gestützte Transformation
heterogener Engineering-Information in SysML~v2.

Kapitel~\ref{sec:implementation} beschreibt anschließend die prototypische
Realisierung dieses Architekturkonzepts im Turing Generator. Der Prototyp
operationalisiert die entwickelten Mechanismen und ermöglicht deren
anschließende Untersuchung anhand konkreter Transformationen.